% ============================================================================
% KVault-DB: A Comprehensive Implementation Guide
% Compile with: pdflatex KVault_Internals.tex (run twice for ToC)
% Tested with: TeX Live 2023, MiKTeX 23.x, Overleaf
% ============================================================================
\documentclass[11pt,a4paper]{article}

% --- Core Packages ---
\usepackage[margin=1in]{geometry}
\usepackage[utf8]{inputenc}
\usepackage[T1]{fontenc}
\usepackage{lmodern}
\usepackage{microtype}

% --- Math ---
\usepackage{amsmath}
\usepackage{amssymb}

% --- Graphics ---
\usepackage{graphicx}

% --- Hyperlinks & PDF Metadata ---
\usepackage[
    colorlinks=true,
    linkcolor=blue!70!black,
    citecolor=green!50!black,
    urlcolor=blue!70!black,
    pdftitle={KVault-DB: A Comprehensive Implementation Guide},
    pdfauthor={Systems Engineering Documentation},
    pdfsubject={C++20 LSM-Tree Key-Value Store},
]{hyperref}

% --- Code Listings ---
\usepackage{listings}
\usepackage{xcolor}

% --- Define colour palette ---
\definecolor{codebg}{RGB}{248,248,248}
\definecolor{codeframe}{RGB}{200,200,200}
\definecolor{kwcolor}{RGB}{0,0,180}
\definecolor{strcolor}{RGB}{163,21,21}
\definecolor{cmtcolor}{RGB}{0,128,0}
\definecolor{numcolor}{RGB}{128,128,128}
\definecolor{prepcolor}{RGB}{128,0,128}

% --- C++ listing style ---
\lstdefinestyle{cpp}{
    language=C++,
    basicstyle=\ttfamily\small,
    keywordstyle=\color{kwcolor}\bfseries,
    stringstyle=\color{strcolor},
    commentstyle=\color{cmtcolor}\itshape,
    directivestyle=\color{prepcolor},
    numberstyle=\tiny\color{numcolor},
    backgroundcolor=\color{codebg},
    rulecolor=\color{codeframe},
    numbers=left,
    numbersep=8pt,
    stepnumber=1,
    frame=single,
    framesep=4pt,
    breaklines=true,
    breakatwhitespace=false,
    showstringspaces=false,
    tabsize=4,
    captionpos=b,
    morekeywords={
        constexpr, noexcept, nullptr, override, final,
        std, unique_ptr, shared_ptr, optional, vector,
        string, string_view, filesystem, shared_mutex,
        shared_lock, unique_lock, uint8_t, uint32_t,
        uint64_t, size_t, int64_t
    },
}

% --- JavaScript listing style ---
\lstdefinestyle{js}{
    basicstyle=\ttfamily\small,
    keywordstyle=\color{kwcolor}\bfseries,
    stringstyle=\color{strcolor},
    commentstyle=\color{cmtcolor}\itshape,
    numberstyle=\tiny\color{numcolor},
    backgroundcolor=\color{codebg},
    rulecolor=\color{codeframe},
    numbers=left,
    numbersep=8pt,
    frame=single,
    breaklines=true,
    showstringspaces=false,
    tabsize=2,
    captionpos=b,
    morekeywords={const,let,import,export,from,async,await,
                  useState,useEffect,useCallback,useRef},
}

% --- CMake listing style ---
\lstdefinestyle{cmake}{
    language=make,
    basicstyle=\ttfamily\small,
    keywordstyle=\color{kwcolor}\bfseries,
    commentstyle=\color{cmtcolor}\itshape,
    numberstyle=\tiny\color{numcolor},
    backgroundcolor=\color{codebg},
    rulecolor=\color{codeframe},
    numbers=left,
    numbersep=8pt,
    frame=single,
    breaklines=true,
    showstringspaces=false,
    captionpos=b,
}

% --- Default style ---
\lstset{style=cpp}

% --- Misc formatting ---
\usepackage{enumitem}
\usepackage{booktabs}
\usepackage{array}
\usepackage{fancyhdr}

\pagestyle{fancy}
\fancyhf{}
\lhead{\textit{KVault-DB: A Comprehensive Implementation Guide}}
\rhead{\thepage}
\setlength{\headheight}{14pt}

% ============================================================================
% TITLE PAGE
% ============================================================================
\title{
    \vspace{-0.5cm}
    \LARGE\textbf{KVault-DB}\\[0.4cm]
    \large A Comprehensive, Step-by-Step Implementation Guide\\
    \large to a Full-Stack C++20 LSM-Tree Key-Value Store\\[0.6cm]
    \normalsize\textit{From shared-memory primitives to a live React dashboard ---}\\
    \normalsize\textit{every layer, every line, explained.}
}
\author{
    Systems Engineering Documentation\\
    \small Based on the KVault-DB open-source codebase
}
\date{\today}

% ============================================================================
\begin{document}
% ============================================================================

\maketitle
\thispagestyle{empty}
\vspace{2cm}

\begin{abstract}
\noindent
KVault-DB is a full-stack, single-node Key-Value Store built from first
principles in modern C++20. Its storage engine implements the
Log-Structured Merge-Tree (LSM-Tree) architecture --- the same foundational
design used in production systems such as LevelDB, RocksDB, and Apache
Cassandra. The system exposes its engine through a multi-threaded HTTP REST API
(built with the Crow microframework) and pairs it with a real-time React/Vite
visualization dashboard.

This document dissects the implementation of KVault-DB into ten logical,
sequential construction steps. Each step covers the theoretical motivation
for the design choice, the practical build process, the key architectural
decisions, and the exact C++ code extracted directly from the source files.
By the end, the reader will have a complete, bottom-up understanding of
how a production-grade storage engine is assembled from raw system primitives.
\end{abstract}

\newpage
\tableofcontents
\newpage

% ============================================================================
\section{Step 1: Foundation --- Core Types and Engine Configuration}
\label{sec:types}
% ============================================================================

\subsection{Conceptual Overview}

Before writing a single line of business logic, a storage engine must
establish a shared vocabulary. Every layer of the system --- from the WAL
serializer to the HTTP API --- must agree on what a \emph{record} is and what
types of mutations are possible. Defining these types centrally prevents
circular dependencies and eliminates the risk of subtle semantic mismatches
between layers.

The most important design choice at this layer is the \textbf{Tombstone
Sentinel}. In an LSM-Tree, data is \emph{never physically deleted} at
insertion time. The storage is append-only. A \texttt{DELETE(``x'')} operation
is not a removal; it is the \emph{insertion of a deletion marker}
(a tombstone). The system then handles the semantics of deletion during reads
and background compaction.

\subsection{The Building Process}

The foundation consists of two files:
\begin{itemize}[noitemsep]
    \item \texttt{include/kvault/types.hpp} --- primitive type aliases and the \texttt{KVRecord} struct.
    \item \texttt{include/kvault/config.hpp} --- the \texttt{EngineConfig} struct with all tuneable parameters.
\end{itemize}

The \texttt{EngineConfig} struct acts as a dependency-injection container. It is
constructed in \texttt{main.cpp} and passed into the \texttt{KVStore} constructor,
which distributes the relevant fields to each subsystem (WAL, MemTable, SSTable
Manager). This design means no subsystem reads global variables, making the
entire engine independently testable.

\subsection{Implementation Details}

The \texttt{RecordType} enum uses \texttt{uint8\_t} as its underlying type.
This is not an accident: the WAL binary serializer writes exactly one byte for
the record type field. Using \texttt{uint8\_t} guarantees that the C++ type
and the binary wire format are identical, eliminating padding or alignment
surprises.

The tombstone sentinel \texttt{kTombstoneValue} is a \texttt{constexpr
string\_view} with sentinel bytes (\texttt{0x7F}) surrounding a unique marker
string. It is stored as the \emph{value} of a deleted key inside the SkipList.
This design avoids a separate boolean field on each node, keeping the node
layout flat. When the MemTable reads back this sentinel, it interprets the
key as deleted and returns \texttt{std::nullopt} to the caller.

\subsection{Code \& Configuration}

\begin{lstlisting}[style=cpp, caption={Core Types: \texttt{include/kvault/types.hpp}}]
#pragma once
#include <cstdint>
#include <string>
#include <string_view>

namespace kvault {

// Both are variable-length strings. Keys are compared lexicographically
// for ordered storage in the Skip List and SSTables.
using Key   = std::string;
using Value = std::string;

// Stored as a single byte in the WAL binary format and SSTable records.
// uint8_t ensures exact 1-byte serialization with no padding.
enum class RecordType : uint8_t {
    PUT    = 0x00,  // Insert or overwrite a key-value pair
    DELETE = 0x01,  // Remove a key (written as a tombstone in the MemTable)
};

// The universal mutation unit. Every write flowing through the system
// is represented as a KVRecord.
struct KVRecord {
    RecordType type;
    Key        key;
    Value      value;   // Empty for DELETE records
};

// Used internally by the MemTable to mark deleted keys.
// 0x7F sentinel bytes make collision with real user data virtually impossible.
// NEVER exposed to external callers.
inline constexpr std::string_view kTombstoneValue =
    "\x7F__KVAULT_TOMBSTONE__\x7F";

} // namespace kvault
\end{lstlisting}

\begin{lstlisting}[style=cpp, caption={Engine Configuration: \texttt{include/kvault/config.hpp} (excerpt)}]
namespace kvault {

struct EngineConfig {
    // Approximate byte threshold at which the MemTable is flushed to disk.
    // Default: 4 MB. Larger values increase write throughput (fewer flushes)
    // but consume more RAM.
    size_t memtable_flush_threshold_bytes = 4UL * 1024 * 1024;

    // Maximum Skip List levels. 16 levels supports up to ~65K entries
    // per MemTable flush cycle.
    int skip_list_max_level = 16;

    // WAL and SSTable storage directories.
    std::string wal_directory     = "data/wal";
    std::string sstable_directory = "data/sstables";

    // If true, every WAL append issues fsync() to guarantee stable storage.
    // If false, faster but risks losing last few records on power failure.
    bool sync_per_write = true;

    // Bits allocated per key in each Bloom Filter.
    // 10 bits/key -> ~0.8% false positive rate (k ~= 7 hash functions).
    size_t bloom_filter_bits_per_key = 10;

    // Port for the embedded Crow HTTP server.
    uint16_t server_port = 8080;
};

} // namespace kvault
\end{lstlisting}

% ============================================================================
\section{Step 2: The Probabilistic In-Memory Index --- The Skip List}
\label{sec:skiplist}
% ============================================================================

\subsection{Conceptual Overview}

The MemTable needs a data structure that keeps keys \emph{sorted} at all times
(so flushing to a sorted SSTable is a simple sequential scan) while supporting
$O(\log n)$ insertions, deletions, and exact-key searches.

A \textbf{Skip List} is a probabilistic ordered data structure invented by
William Pugh (1990). It is a layered linked list where each node participates
in one or more \emph{express lanes} at higher levels. A search starts at the
highest level, skipping large chunks of the list, then descends to lower levels
as it narrows in on the target key. The expected complexity for all operations
is $O(\log n)$, equivalent to a balanced binary search tree but with a far
simpler implementation that is amenable to concurrent access.

A node of height $h$ is promoted to level $h+1$ with probability $p = 0.5$.
This gives an expected node height of $\frac{1}{1-p} = 2$, meaning the
expected total number of forward pointers across all $n$ nodes is $2n$ ---
highly space-efficient.

\subsection{The Building Process}

A na\"ive implementation allocates each node with \texttt{new} and deletes
them recursively. For large lists (millions of nodes), the recursive
destructor causes a \textbf{stack overflow}. KVault solves this with an
\textbf{Arena Allocator}: every node is stored in a
\texttt{std::vector<std::unique\_ptr<SkipListNode>>} called \texttt{arena\_}.
When the list is destroyed, the vector's destructor calls \texttt{delete} on
each node iteratively, not recursively.

\subsection{Implementation Details}

The critical insight is the \textbf{ownership split}:
\begin{itemize}[noitemsep]
    \item \texttt{arena\_} holds \texttt{unique\_ptr<Node>} --- it \emph{owns} every node.
    \item \texttt{forward[]} inside each node holds raw \texttt{Node*} pointers --- non-owning, for traversal only.
\end{itemize}

Using \texttt{unique\_ptr} for \texttt{forward[]} would be incorrect because
the same node is reachable from multiple levels simultaneously. Two
\texttt{unique\_ptr}s pointing to the same node would cause a double-free
Undefined Behaviour.

The \texttt{insert()} algorithm works by building an \texttt{update[]}
vector recording the last node at each level whose key is less than the
target. After the traversal, a new node is spliced between
\texttt{update[i]} and \texttt{update[i]->forward[i]} at each level.

\subsection{Code \& Configuration}

\begin{lstlisting}[style=cpp, caption={Skip List Node and Class Declaration: \texttt{include/kvault/skiplist.hpp}}]
struct SkipListNode {
    Key   key;
    Value value;

    // Non-owning traversal pointers. Lifetime managed by SkipList::arena_.
    // Using raw pointers here is correct: the same node can be pointed to
    // from multiple levels, so unique_ptr would cause a double-free.
    std::vector<SkipListNode*> forward;

    // Data node constructor
    SkipListNode(Key k, Value v, int height)
        : key(std::move(k))
        , value(std::move(v))
        , forward(static_cast<size_t>(height), nullptr) {}

    // Sentinel (head) node constructor
    explicit SkipListNode(int height)
        : forward(static_cast<size_t>(height), nullptr) {}
};

class SkipList {
public:
    static constexpr int    kMaxLevel    = 16;
    static constexpr double kProbability = 0.5;

    // Destructor: arena_ iteratively destroys all nodes (O(n), stack-safe).
    ~SkipList() = default;

    void                  insert(const Key& key, const Value& value);
    std::optional<Value>  search(const Key& key) const;
    bool                  remove(const Key& key);

private:
    int random_level();

    std::unique_ptr<SkipListNode> head_;   // Sentinel, entry point at all levels
    std::vector<std::unique_ptr<SkipListNode>> arena_; // OWNS all data nodes
    int    current_level_;
    size_t size_;
    mutable std::mt19937 rng_;
    mutable std::uniform_real_distribution<double> dist_;
};
\end{lstlisting}

\begin{lstlisting}[style=cpp, caption={Skip List Insert and Search: \texttt{src/skiplist.cpp}}]
int SkipList::random_level() {
    int level = 1;
    while (level < kMaxLevel && dist_(rng_) < kProbability)
        ++level;
    return level;
}

void SkipList::insert(const Key& key, const Value& value) {
    // update[i] = rightmost node at level i whose key < target.
    std::vector<SkipListNode*> update(kMaxLevel, nullptr);
    SkipListNode* current = head_.get();

    // Traverse: start at highest level, descend when key would overshoot.
    for (int i = current_level_ - 1; i >= 0; --i) {
        while (current->forward[i] && current->forward[i]->key < key)
            current = current->forward[i];
        update[i] = current;
    }

    // Upsert: if key exists at level 0, overwrite in place.
    SkipListNode* candidate = current->forward[0];
    if (candidate && candidate->key == key) {
        candidate->value = value;
        return;
    }

    // Allocate new node in the arena (unique ownership).
    const int new_height = random_level();
    if (new_height > current_level_) {
        for (int i = current_level_; i < new_height; ++i)
            update[i] = head_.get();
        current_level_ = new_height;
    }

    arena_.push_back(std::make_unique<SkipListNode>(key, value, new_height));
    SkipListNode* new_node = arena_.back().get(); // raw ptr for wiring

    // Splice new_node between update[i] and update[i]->forward[i].
    for (int i = 0; i < new_height; ++i) {
        new_node->forward[i]    = update[i]->forward[i];
        update[i]->forward[i]   = new_node;
    }
    ++size_;
}

std::optional<Value> SkipList::search(const Key& key) const {
    const SkipListNode* current = head_.get();
    for (int i = current_level_ - 1; i >= 0; --i)
        while (current->forward[i] && current->forward[i]->key < key)
            current = current->forward[i];

    const SkipListNode* candidate = current->forward[0];
    if (candidate && candidate->key == key)
        return candidate->value;
    return std::nullopt;
}
\end{lstlisting}

% ============================================================================
\section{Step 3: The In-Memory Write Buffer --- The MemTable}
\label{sec:memtable}
% ============================================================================

\subsection{Conceptual Overview}

The \textbf{MemTable} is the active, mutable tier of the LSM-Tree. Every
incoming \texttt{PUT} and \texttt{DELETE} request is directed here first. It
acts as a write buffer: mutations accumulate in memory until the table reaches
a configurable byte threshold (e.g., 4~MB), at which point it is
\emph{frozen} (made immutable) and flushed to disk as a new SSTable file.

This design transforms the write problem from a random-write workload (the
bane of mechanical hard drives and flash storage) into a sequential write
workload, dramatically improving throughput. The trade-off is read amplification:
a lookup may need to search the MemTable and multiple SSTable files on disk.

\subsection{The Building Process}

The \texttt{MemTable} class wraps the \texttt{SkipList} and adds:
\begin{enumerate}[noitemsep]
    \item \textbf{Byte tracking}: it maintains a running \texttt{approximate\_size\_} counter to know when to trigger a flush.
    \item \textbf{Tombstone logic}: \texttt{remove()} inserts the key with \texttt{kTombstoneValue} as its value.
    \item \textbf{Snapshot generation}: \texttt{snapshot()} walks the Skip List's sorted level-0 chain and emits a \texttt{std::vector<KVRecord>} for the SSTable writer.
\end{enumerate}

\subsection{Implementation Details}

The tombstone design is central to LSM-Tree correctness. When a user calls
\texttt{DELETE(``x'')}, the key might exist in an older SSTable on disk.
If the MemTable simply performed a physical removal from the Skip List, the
subsequent read path would find nothing in the MemTable and then \emph{fall
through to the SSTable and return stale data}. By inserting the tombstone,
the MemTable short-circuits this: \texttt{contains\_tombstone(key)} returns
\texttt{true}, and the \texttt{KVStore} never initiates the expensive disk search.

\subsection{Code \& Configuration}

\begin{lstlisting}[style=cpp, caption={MemTable Core Logic: \texttt{src/memtable.cpp}}]
void MemTable::put(const Key& key, const Value& value) {
    auto existing = skiplist_.search(key);
    if (existing.has_value()) {
        // Upsert: node overhead unchanged; only value size delta matters.
        approximate_size_ -= existing->size();
        approximate_size_ += value.size();
    } else {
        // New entry: account for key, value, and node pointer overhead.
        approximate_size_ += key.size() + value.size() + kNodeOverheadBytes;
    }
    skiplist_.insert(key, value);
}

// Deletes are NOT physical removals. We insert a tombstone sentinel so
// the read path can short-circuit SSTable lookups for deleted keys.
void MemTable::remove(const Key& key) {
    auto existing = skiplist_.search(key);
    if (existing.has_value()) {
        approximate_size_ -= existing->size();
        approximate_size_ += kTombstoneValue.size();
    } else {
        approximate_size_ += key.size() + kTombstoneValue.size()
                           + kNodeOverheadBytes;
    }
    skiplist_.insert(key, std::string(kTombstoneValue));
}

std::optional<Value> MemTable::get(const Key& key) const {
    auto result = skiplist_.search(key);
    // A tombstone means the key was deleted; return nullopt to the caller.
    if (result.has_value() && *result == kTombstoneValue)
        return std::nullopt;
    return result;
}

// Used by KVStore to short-circuit expensive SSTable lookups for deleted keys.
bool MemTable::contains_tombstone(const Key& key) const {
    auto result = skiplist_.search(key);
    return result.has_value() && *result == kTombstoneValue;
}

// Walk the sorted level-0 chain; produce a sorted KVRecord vector for the
// SSTable writer. This is the 'snapshot' taken before a flush.
std::vector<KVRecord> MemTable::snapshot() const {
    std::vector<KVRecord> records;
    records.reserve(skiplist_.size());
    for (auto it = skiplist_.begin(); it != skiplist_.end(); ++it) {
        KVRecord record;
        if (it.value() == kTombstoneValue) {
            record.type = RecordType::DELETE;
            record.key  = it.key();
        } else {
            record.type  = RecordType::PUT;
            record.key   = it.key();
            record.value = it.value();
        }
        records.push_back(std::move(record));
    }
    return records;
}
\end{lstlisting}

% ============================================================================
\section{Step 4: Crash Resilience --- The Write-Ahead Log}
\label{sec:wal}
% ============================================================================

\subsection{Conceptual Overview}

The MemTable lives in volatile RAM. A server crash --- a power failure, an
OS panic, or even a \texttt{kill -9} signal --- would erase all
unacknowledged writes. The \textbf{Write-Ahead Log (WAL)} is the mechanism that
prevents this data loss.

The WAL is an append-only binary file. The \emph{durability contract} is strict:
\begin{enumerate}[noitemsep]
    \item A write request arrives at the \texttt{KVStore}.
    \item The mutation is serialized and appended to the WAL file on disk.
    \item \textbf{Only after} the WAL write is confirmed does the MemTable get updated.
    \item The HTTP response is sent to the client.
\end{enumerate}
On startup, the WAL is \emph{replayed}: every record is re-applied to a fresh
MemTable, reconstructing the exact in-memory state at the time of the crash.

\subsection{The Building Process}

The WAL is a self-describing binary log. Each record has a fixed-layout header
and a trailing CRC32 checksum. This allows the replay reader to detect
\emph{partial writes} --- a crash that interrupted the write mid-record leaves
a corrupted final record. The replay stops at the first CRC mismatch, safely
discarding all records after it.

The CRC32 lookup table is generated at \textbf{compile time} using
\texttt{constexpr}. At runtime, computing a checksum is a simple table lookup
with no initialization overhead.

\subsection{Implementation Details}

The durability is achieved via two levels of flushing:
\begin{itemize}[noitemsep]
    \item \texttt{fflush()} --- pushes data from the C library's userspace buffer into the OS kernel's page cache (fast, in RAM).
    \item \texttt{fsync()} / \texttt{\_commit()} --- forces the OS to write the page cache to the physical storage device (slow, but survives a power failure).
\end{itemize}
When \texttt{sync\_per\_write = true} (the default), both are called on every
append. This limits throughput to the hardware's IOPS capacity (typically
500--5000 per second for consumer SSDs) but provides 100\% crash safety.

\subsection{Code \& Configuration}

\begin{lstlisting}[style=cpp, caption={Compile-Time CRC32 Table and WAL Binary Format: \texttt{src/wal.cpp}}]
// Standard CRC32, reflected polynomial 0xEDB88320 (same as Ethernet/zlib).
// Generated entirely at compile time -- zero runtime initialization cost.
constexpr std::array<uint32_t, 256> generate_crc32_table() {
    std::array<uint32_t, 256> table{};
    for (uint32_t i = 0; i < 256; ++i) {
        uint32_t crc = i;
        for (int j = 0; j < 8; ++j)
            crc = (crc & 1u) ? (crc >> 1) ^ 0xEDB88320u : (crc >> 1);
        table[i] = crc;
    }
    return table;
}
constexpr auto kCRC32Table = generate_crc32_table();

// WAL Binary Record Layout (all multi-byte fields are little-endian):
//
// Offset  Size  Field         Description
// ------  ----  ----------    -------------------------------------------
// 0       1     record_type   0x00 = PUT, 0x01 = DELETE
// 1       4     key_length    Length of key_data in bytes (uint32 LE)
// 5       K     key_data      Raw key bytes
// 5+K     4     value_length  Length of value_data in bytes (uint32 LE)
// 9+K     V     value_data    Raw value bytes
// 9+K+V   4     crc32         CRC32 over all preceding bytes (uint32 LE)
//                             Total: 13 + K + V bytes per record.

std::vector<uint8_t> WriteAheadLog::serialize_record(const KVRecord& record) {
    const auto key_len = static_cast<uint32_t>(record.key.size());
    const auto val_len = static_cast<uint32_t>(record.value.size());
    std::vector<uint8_t> buf;
    buf.reserve(13 + key_len + val_len);

    buf.push_back(static_cast<uint8_t>(record.type)); // 1 byte type
    write_u32_le(buf, key_len);                        // 4 bytes key length
    buf.insert(buf.end(), record.key.begin(), record.key.end());
    write_u32_le(buf, val_len);                        // 4 bytes value length
    buf.insert(buf.end(), record.value.begin(), record.value.end());

    // Append CRC32 of all preceding bytes to detect partial writes.
    const uint32_t crc = compute_crc32(buf.data(), buf.size());
    write_u32_le(buf, crc);
    return buf;
}
\end{lstlisting}

\begin{lstlisting}[style=cpp, caption={WAL Append and Platform Sync: \texttt{src/wal.cpp}}]
void WriteAheadLog::append(const KVRecord& record) {
    const auto serialized = serialize_record(record);
    const size_t written = std::fwrite(
        serialized.data(), 1, serialized.size(), write_handle_);
    if (written != serialized.size())
        throw std::runtime_error("WAL: short write");
    sync_to_disk();
}

void WriteAheadLog::sync_to_disk() {
    if (!write_handle_) return;
    std::fflush(write_handle_);    // Push userspace buffer to OS page cache.
    if (sync_per_write_) {
#ifdef _WIN32
        _commit(_fileno(write_handle_));   // Windows hardware flush
#else
        ::fsync(fileno(write_handle_));    // POSIX hardware flush
#endif
    }
}

// Replay: re-applies all valid WAL records to a fresh MemTable on startup.
std::vector<KVRecord> WriteAheadLog::replay() const {
    std::vector<KVRecord> records;
    FILE* read_handle = std::fopen(path_.string().c_str(), "rb");
    if (!read_handle) return records; // First run, no WAL yet.
    while (true) {
        auto record = deserialize_record(read_handle);
        if (!record.has_value()) break; // EOF or CRC mismatch -- stop here.
        records.push_back(std::move(*record));
    }
    std::fclose(read_handle);
    return records;
}
\end{lstlisting}

% ============================================================================
\section{Step 5: Read Optimization --- The Bloom Filter}
\label{sec:bloom}
% ============================================================================

\subsection{Conceptual Overview}

In an LSM-Tree, a \texttt{GET} request that misses the MemTable must search
the SSTable files on disk. Disk I/O is orders of magnitude slower than RAM
access. A key that has \emph{never been inserted} into the database would
cause the system to search every SSTable file before concluding the key does
not exist --- a needlessly expensive operation.

A \textbf{Bloom Filter} is a space-efficient probabilistic data structure that
answers the set membership question. Given a set $S$ and a query element $x$:
\begin{itemize}[noitemsep]
    \item \textbf{FALSE} $\Rightarrow$ $x$ is \emph{definitively} not in $S$.
    \item \textbf{TRUE} $\Rightarrow$ $x$ is \emph{probably} in $S$ (false positives possible).
\end{itemize}
False negatives are mathematically impossible. This asymmetry makes Bloom
Filters a perfect pre-filter for disk I/O: a \texttt{FALSE} result
eliminates the SSTable search entirely.

With $m$ bits and $k$ hash functions for $n$ keys, the false positive rate is:
\[ p = \left(1 - e^{-kn/m}\right)^k \]
Optimal values: $m = -n \ln(p) / (\ln 2)^2$, $k = (m/n) \ln 2$.

With \texttt{bits\_per\_key = 10}: $p \approx 0.0082$ ($\approx 0.8\%$ false positive rate).

\subsection{The Building Process}

Instead of implementing $k$ independent hash functions (complex and expensive),
KVault uses the \textbf{Kirsch--Mitzenmacher double hashing technique} [ESA 2006].
Two base hashes $h_1$ and $h_2$ are computed from the key using FNV-1a and a
Murmur-inspired finalizer. The $k$ composite hashes are then:
\[g_i(\text{key}) = (h_1(\text{key}) + i \cdot h_2(\text{key})) \bmod m, \quad i \in [0, k)\]

\subsection{Implementation Details}

$h_2$ is forced to be odd via \texttt{h2 |= 1}. This guarantees that
$h_2$ is coprime with any power-of-two $m$, ensuring the sequence
$h_1, h_1+h_2, h_1+2h_2, \ldots$ visits $m$ distinct positions before
repeating (full coverage of the bit array).

The filter is serialized as \texttt{[4 bytes: k][4 bytes: m][ceil(m/8) bytes: bits]}
and embedded directly into the SSTable file footer region, so it can be loaded
into memory without a separate I/O call.

\subsection{Code \& Configuration}

\begin{lstlisting}[style=cpp, caption={Bloom Filter Double Hashing: \texttt{src/bloom\_filter.cpp}}]
// Two independent 64-bit hashes from the same key.
// h1: FNV-1a -- fast, excellent avalanche for short keys.
// h2: Murmur finalizer applied to h1 -- statistically independent.
std::pair<uint64_t, uint64_t> BloomFilter::base_hashes(const Key& key) {
    constexpr uint64_t kFNVPrime  = 0x00000100000001B3ULL;
    constexpr uint64_t kFNVOffset = 0xCBF29CE484222325ULL;

    uint64_t h1 = kFNVOffset;
    for (char ch : key) {
        h1 ^= static_cast<uint64_t>(static_cast<unsigned char>(ch));
        h1 *= kFNVPrime;
    }

    // Murmur finalizer mix: breaks linear correlations in h1.
    uint64_t h2 = h1;
    h2 ^= h2 >> 33; h2 *= 0xFF51AFD7ED558CCDULL;
    h2 ^= h2 >> 33; h2 *= 0xC4CEB9FE1A85EC53ULL;
    h2 ^= h2 >> 33;
    return { h1, h2 };
}

void BloomFilter::add(const Key& key) {
    auto [h1, h2] = base_hashes(key);
    h2 |= 1; // Ensure h2 is odd (coprime with m for full bit-array coverage)

    for (size_t i = 0; i < k_; ++i) {
        const size_t bit_pos = (h1 + i * h2) % m_;
        bits_[bit_pos / 8] |= static_cast<uint8_t>(1u << (bit_pos % 8));
    }
}

bool BloomFilter::might_contain(const Key& key) const {
    auto [h1, h2] = base_hashes(key);
    h2 |= 1;

    for (size_t i = 0; i < k_; ++i) {
        const size_t bit_pos = (h1 + i * h2) % m_;
        if (!(bits_[bit_pos / 8] & (1u << (bit_pos % 8))))
            return false; // Definite miss -- key was never add()ed
    }
    return true; // Probable hit -- proceed to disk search
}
\end{lstlisting}

% ============================================================================
\section{Step 6: Disk Persistence --- The SSTable Writer}
\label{sec:sstable_writer}
% ============================================================================

\subsection{Conceptual Overview}

When the MemTable reaches its size threshold, it must be written to disk as an
immutable \textbf{Sorted String Table (SSTable)}. Because the MemTable's Skip
List is already sorted by key, the flush is a simple, single-pass sequential
write --- one of the most efficient I/O patterns available on any storage medium.

An SSTable is a self-describing binary file. A reader starting from any state
can load an SSTable, locate its Bloom Filter, and binary-search for any key
without any external metadata.

\subsection{The Building Process}

The SSTable file layout is written sequentially in four phases:
\begin{enumerate}[noitemsep]
    \item \textbf{Data Block}: all key-value records encoded sequentially.
    \item \textbf{Sparse Index Block}: one index entry per 100 keys (configurable), each storing a key and its byte offset in the data block. Enables $O(\log n)$ block-level lookup.
    \item \textbf{Bloom Filter Block}: the serialized Bloom Filter (enables $O(1)$ membership test before any data block I/O).
    \item \textbf{Footer} (fixed 48 bytes): byte offsets of all three blocks, entry count, and a magic number for file integrity validation.
\end{enumerate}

\subsection{Implementation Details}

The \texttt{SSTableFooter} struct uses \texttt{static\_assert} to enforce its
exact byte size at compile time:

\begin{equation*}
    \texttt{static\_assert(sizeof(SSTableFooter) == 48)}
\end{equation*}

This is a zero-cost compile-time contract. If a developer ever accidentally
adds a new field and forgets to update serialization code, the build fails
immediately rather than silently producing a corrupted SSTable at runtime.

\subsection{Code \& Configuration}

\begin{lstlisting}[style=cpp, caption={SSTable File Layout and Footer: \texttt{include/kvault/sstable\_writer.hpp}}]
// SSTable Binary Layout:
//
//  +---------------------------+  <- byte 0
//  |  DATA BLOCK               |
//  |  Sequential KVRecords     |
//  +---------------------------+  <- index_block_offset
//  |  SPARSE INDEX BLOCK       |
//  |  [key_len][key][offset]   |  (every kIndexBlockInterval-th key)
//  +---------------------------+  <- bloom_block_offset
//  |  BLOOM FILTER BLOCK       |
//  |  [k][m][bit_array...]     |
//  +---------------------------+  <- footer_offset
//  |  FOOTER (48 bytes)        |
//  |  Block offsets + magic    |
//  +---------------------------+  <- EOF

class SSTableWriter {
public:
    // One sparse index entry per 100 records. The reader binary-searches
    // this index, then linearly scans within the block.
    static constexpr size_t kIndexBlockInterval = 100;

    // Magic number at end of footer: ASCII "KVLT" + version 0x0001.
    static constexpr uint64_t kMagicNumber = 0x544C564B00010000ULL;

    // Write a sorted vector of KVRecords to a new SSTable file.
    // sorted_entries MUST be in ascending key order (from MemTable::snapshot()).
    static void write(const std::filesystem::path& path,
                      const std::vector<KVRecord>& sorted_entries,
                      size_t bits_per_key = 10);

    SSTableWriter() = delete; // Static-only utility class.
};

// Fixed 48-byte footer. static_assert enforces layout at compile time.
struct SSTableFooter {
    uint64_t index_block_offset;  // Byte offset of the Sparse Index Block
    uint64_t bloom_block_offset;  // Byte offset of the Bloom Filter Block
    uint64_t footer_offset;       // Byte offset of this footer itself
    uint64_t entry_count;         // Total number of records in the SSTable
    uint64_t data_block_size;     // Size of the Data Block in bytes
    uint64_t magic;               // kMagicNumber -- validates file integrity
    static constexpr size_t kSerializedSize = 48;
};
// Compile-time contract: if a field is added, the build fails.
static_assert(sizeof(SSTableFooter) == SSTableFooter::kSerializedSize,
              "SSTableFooter size mismatch -- update kSerializedSize");
\end{lstlisting}

% ============================================================================
\section{Step 7: Reading from Disk --- The SSTable Reader}
\label{sec:sstable_reader}
% ============================================================================

\subsection{Conceptual Overview}

Reading from an SSTable is the inverse of writing. The reader must:
\begin{enumerate}[noitemsep]
    \item Load the 48-byte footer to locate the Bloom Filter and Sparse Index.
    \item Check the Bloom Filter. If it returns \texttt{FALSE}, return immediately.
    \item Binary-search the Sparse Index to find the data block offset nearest to the target key.
    \item Linearly scan the data block from that offset until the key is found or passed.
\end{enumerate}

This three-tier search strategy minimizes disk I/O. The Bloom Filter lives in
RAM. The Sparse Index is small enough to keep in RAM. Only the final linear
scan touches the potentially large data block on disk.

\subsection{Implementation Details}

On construction, the \texttt{SSTableReader} reads the footer and loads the
Bloom Filter into memory. This is a \emph{one-time cost} paid when the
SSTable Manager loads files at startup.

Lexicographic key ordering, which the Skip List and SSTable Writer both
preserve, is the key invariant that makes binary search on the Sparse Index
correct. In C++, \texttt{std::string} comparison is lexicographic by default,
so no custom comparator is needed.

\subsection{Code \& Configuration}

\begin{lstlisting}[style=cpp, caption={SSTable Search Strategy (conceptual from \texttt{src/sstable\_reader.cpp})}]
// On construction: load footer + Bloom Filter into RAM (one-time cost).
// SSTableReader::SSTableReader(const std::filesystem::path& path) {
//   1. Open file.
//   2. Seek to EOF - 48 bytes, read footer.
//   3. Validate magic number.
//   4. Seek to footer.bloom_block_offset, read Bloom Filter bytes.
//   5. Deserialize BloomFilter from bytes. Store in bloom_filter_.
// }

// get_record(): The three-tier lookup.
// std::optional<KVRecord> SSTableReader::get_record(const Key& key) const {
//
//   // Tier 1: Bloom Filter (in RAM, O(k) time, no disk I/O)
//   if (!bloom_filter_.might_contain(key))
//       return std::nullopt; // Definite miss. Zero disk I/O needed.
//
//   // Tier 2: Binary search the Sparse Index (small, usually cached by OS)
//   // to find the data block offset nearest to 'key'.
//   uint64_t scan_offset = binary_search_index(key);
//
//   // Tier 3: Linear scan the data block from scan_offset.
//   // Decode records sequentially until key matches or key is passed.
//   seek_to(scan_offset);
//   while (not_at_index_block) {
//       KVRecord rec = decode_next_record();
//       if (rec.key == key) return rec;  // Found
//       if (rec.key >  key) return std::nullopt; // Passed it -- not present
//   }
//   return std::nullopt;
// }
\end{lstlisting}

% ============================================================================
\section{Step 8: Managing Generations --- The SSTable Manager}
\label{sec:sstable_manager}
% ============================================================================

\subsection{Conceptual Overview}

As the engine runs, each MemTable flush produces a new \texttt{.sst} file on
disk. Over time, there may be dozens or hundreds of them. The
\textbf{SSTable Manager} is responsible for keeping track of all active SSTable
files and executing the \emph{cascading read path}: searching SSTables from
newest to oldest until the key is found.

The generation ordering is critical. If key \texttt{``x''} was written with
value \texttt{``v1''} and later overwritten with \texttt{``v2''}, these may
reside in different SSTables. Searching newest-to-oldest ensures that
\texttt{``v2''} (the newer, correct value) is returned without ever reading
the stale \texttt{``v1''} record.

\subsection{Implementation Details}

SSTable filenames are zero-padded millisecond timestamps
(e.g., \texttt{0001691234567000.sst}). This ensures that lexicographic
filename ordering matches chronological ordering. On startup, the Manager sorts
files in \emph{descending} lexicographic order (newest first) so that
\texttt{readers\_[0]} is always the most recent SSTable.

When a new SSTable is flushed, it is prepended to \texttt{readers\_} at
index 0 --- maintaining the newest-first invariant without re-sorting.

\subsection{Code \& Configuration}

\begin{lstlisting}[style=cpp, caption={SSTable Manager: Startup Scan and Cascading Read: \texttt{src/sstable\_manager.cpp}}]
SSTableManager::SSTableManager(const std::filesystem::path& sstable_dir)
    : sstable_dir_(sstable_dir)
{
    std::filesystem::create_directories(sstable_dir_);

    // Collect all .sst files in the directory.
    std::vector<std::filesystem::path> sst_files;
    for (const auto& entry : std::filesystem::directory_iterator(sstable_dir_)) {
        if (entry.is_regular_file() && entry.path().extension() == ".sst")
            sst_files.push_back(entry.path());
    }

    // Sort descending (newest first) by lexicographic filename.
    // Zero-padded timestamps ensure this is also chronological order.
    std::sort(sst_files.begin(), sst_files.end(), std::greater<>());

    for (const auto& path : sst_files)
        readers_.push_back(std::make_unique<SSTableReader>(path));
}

// Cascading read: search newest to oldest.
// The first match wins -- it represents the most recent state of the key.
std::optional<KVRecord> SSTableManager::get_record(const Key& key) const {
    for (const auto& reader : readers_) {
        auto result = reader->get_record(key);
        if (result) {
            // Found! Could be PUT or DELETE (tombstone from an older flush).
            return result;
        }
    }
    return std::nullopt;
}

// A newly flushed SSTable is always the newest generation.
// Prepend to readers_ to maintain newest-first ordering.
void SSTableManager::add_sstable(const std::filesystem::path& path) {
    auto reader = std::make_unique<SSTableReader>(path);
    readers_.insert(readers_.begin(), std::move(reader));
}
\end{lstlisting}

% ============================================================================
\section{Step 9: The Concurrent Engine Core --- KVStore}
\label{sec:kvstore}
% ============================================================================

\subsection{Conceptual Overview}

The \textbf{KVStore} class is the central orchestrator of the entire storage
engine. It owns and coordinates the MemTable, the WAL, and the SSTable Manager
behind a single, thread-safe public interface. Because the Crow HTTP server
handles many client requests concurrently, the \texttt{KVStore} must be able
to safely serve multiple readers and writers simultaneously.

The concurrency strategy is a classic \textbf{Reader-Writer lock}
(\texttt{std::shared\_mutex}):
\begin{itemize}[noitemsep]
    \item \textbf{Readers} (\texttt{get}) acquire a \texttt{std::shared\_lock} --- multiple readers can proceed simultaneously.
    \item \textbf{Writers} (\texttt{put}, \texttt{remove}, flush) acquire a \texttt{std::unique\_lock} --- exclusive access, blocking all readers and other writers.
\end{itemize}

\subsection{Implementation Details}

The \texttt{remove()} method makes a critically correct design decision. In a
naive implementation, one might first call \texttt{get()} to check if the key
exists before deleting it. This would deadlock: \texttt{get()} acquires a
\texttt{shared\_lock}, but \texttt{remove()} already holds a
\texttt{unique\_lock}, and a single thread cannot hold both simultaneously in
\texttt{std::shared\_mutex}.

The solution is a \emph{blind write}: the tombstone is written unconditionally,
regardless of whether the key exists. This is semantically correct in an
LSM-Tree (a DELETE in an empty set is a no-op on replay) and eliminates the
deadlock. This is the same approach used by all production LSM-Tree databases.

The flush pipeline (\texttt{trigger\_flush()}) executes under the
\texttt{unique\_lock} held by the calling write operation. It snapshots the
MemTable, writes the SSTable to disk, registers the new file with the Manager,
clears the MemTable, and truncates the WAL.

\subsection{Code \& Configuration}

\begin{lstlisting}[style=cpp, caption={Thread-Safe Write and Read Paths: \texttt{src/kvstore.cpp}}]
void KVStore::put(const Key& key, const Value& value) {
    std::unique_lock lock(rw_mutex_); // Exclusive: blocks all readers/writers

    // WAL first -- data is durable on disk before MemTable is updated.
    wal_->append(KVRecord{RecordType::PUT, key, value});
    memtable_->put(key, value);

    if (memtable_->should_flush() ||
        wal_->file_size() >= config_.memtable_flush_threshold_bytes)
        trigger_flush();
}

std::optional<Value> KVStore::get(const Key& key) const {
    std::shared_lock lock(rw_mutex_); // Shared: multiple readers allowed

    // Step 1: Tombstone short-circuit -- avoids expensive disk search.
    if (memtable_->contains_tombstone(key))
        return std::nullopt;

    // Step 2: Search active MemTable.
    auto mem_val = memtable_->get(key);
    if (mem_val) return mem_val;

    // Step 3: Fall back to SSTable Manager (newest-to-oldest disk search).
    auto sst_val = sstable_manager_->get(key);
    if (sst_val && sst_val->empty()) return std::nullopt; // SSTable tombstone
    return sst_val;
}

bool KVStore::remove(const Key& key) {
    std::unique_lock lock(rw_mutex_);

    // BLIND WRITE: write tombstone unconditionally.
    // Do NOT call get() here -- that would require shared_lock, causing
    // a deadlock since we already hold unique_lock.
    wal_->append(KVRecord{RecordType::DELETE, key, ""});
    memtable_->remove(key);

    if (memtable_->should_flush() ||
        wal_->file_size() >= config_.memtable_flush_threshold_bytes)
        trigger_flush();
    return true;
}

// Flush pipeline: snapshot -> write SSTable -> register -> reset.
void KVStore::trigger_flush() {
    // Caller MUST hold unique_lock on rw_mutex_.
    auto records = memtable_->snapshot();
    if (records.empty()) return;

    auto sst_path = std::filesystem::path(config_.sstable_directory)
                    / generate_sstable_filename();
    SSTableWriter::write(sst_path, records, config_.bloom_filter_bits_per_key);
    sstable_manager_->add_sstable(sst_path);

    // Reset MemTable and truncate WAL (records now safely in SSTable).
    memtable_ = std::make_unique<MemTable>(config_.memtable_flush_threshold_bytes);
    wal_->truncate();
}
\end{lstlisting}

% ============================================================================
\section{Step 10: Full-Stack Integration --- REST API and React Dashboard}
\label{sec:api_dashboard}
% ============================================================================

\subsection{Conceptual Overview}

The final step connects the C++ storage engine to the outside world via a
standard HTTP REST API and pairs it with a real-time React visualization
dashboard. The API transforms KVault from an embedded library into a
\emph{network database service}: any application in any language can store
and retrieve data using standard HTTP requests.

The Crow C++ microframework handles HTTP parsing, routing, and thread pool
management. Because the \texttt{KVStore} is already thread-safe via its
internal \texttt{shared\_mutex}, the multi-threaded Crow server can handle
concurrent requests without additional synchronization at the API layer.

\subsection{The Building Process}

The API exposes five endpoints:
\begin{itemize}[noitemsep]
    \item \texttt{GET /api/kv/<key>} --- retrieve a value.
    \item \texttt{POST /api/kv} --- insert or update a key (JSON body: \texttt{\{"key": "...", "value": "..."\}}).
    \item \texttt{DELETE /api/kv/<key>} --- delete a key (idempotent tombstone write).
    \item \texttt{GET /api/metrics} --- live engine telemetry (MemTable size, WAL size, SSTable count).
    \item \texttt{GET /api/memtable/snapshot} --- full sorted dump of the active MemTable, consumed by the React SkipList visualizer.
\end{itemize}

The React dashboard uses two custom hooks: \texttt{useKVStore.js} for
CRUD operations and \texttt{useMetrics.js} for polling
\texttt{/api/metrics} every 2 seconds to update the live charts.

\subsection{Implementation Details}

The \texttt{DELETE} macro collision on Windows is handled explicitly
(\texttt{\#undef DELETE}) before the delete route is defined, preventing the
Windows SDK's \texttt{DELETE} macro from interfering with Crow's method-matching DSL.

The server entry point in \texttt{main.cpp} initializes the engine with
a configuration struct, replays the WAL (recovering from any previous crash),
and starts the HTTP server. The entire startup sequence takes under 100~ms
on typical hardware.

\subsection{Code \& Configuration}

\begin{lstlisting}[style=cpp, caption={REST API Route Definitions: \texttt{src/api\_routes.cpp}}]
void ApiServer::setup_routes() {

    // GET /api/kv/<key>
    CROW_ROUTE(app_, "/api/kv/<string>").methods(crow::HTTPMethod::GET)(
        [this](const std::string& key) {
            auto val = store_->get(key);
            if (!val) return crow::response(404, "Key not found");
            return crow::response(200, *val);
        });

    // POST /api/kv  -- body: {"key": "...", "value": "..."}
    CROW_ROUTE(app_, "/api/kv").methods(crow::HTTPMethod::POST)(
        [this](const crow::request& req) {
            auto x = crow::json::load(req.body);
            if (!x || !x.has("key") || !x.has("value"))
                return crow::response(400, "Invalid JSON or missing fields");
            store_->put(std::string(x["key"].s()), std::string(x["value"].s()));
            return crow::response(200, "OK");
        });

    // DELETE /api/kv/<key>
    // Note: #undef DELETE required on Windows to suppress the SDK macro.
#ifdef DELETE
#undef DELETE
#endif
    CROW_ROUTE(app_, "/api/kv/<string>").methods(crow::HTTPMethod::DELETE)(
        [this](const std::string& key) {
            store_->remove(key);
            return crow::response(200, "OK"); // Always 200 (idempotent)
        });

    // GET /api/metrics  -- live engine telemetry
    CROW_ROUTE(app_, "/api/metrics").methods(crow::HTTPMethod::GET)(
        [this]() {
            crow::json::wvalue metrics;
            metrics["memtable_size_bytes"] = store_->memtable_size();
            metrics["wal_size_bytes"]      = store_->wal_size();
            metrics["sstable_count"]       = store_->sstable_count();
            return crow::response(metrics);
        });

    // GET /api/memtable/snapshot  -- sorted MemTable dump for visualizer
    CROW_ROUTE(app_, "/api/memtable/snapshot").methods(crow::HTTPMethod::GET)(
        [this]() {
            auto records = store_->memtable_snapshot();
            crow::json::wvalue::list arr;
            for (const auto& rec : records) {
                crow::json::wvalue entry;
                entry["key"]   = rec.key;
                entry["value"] = rec.value;
                entry["type"]  = (rec.type == RecordType::DELETE)
                                 ? "tombstone" : "put";
                arr.push_back(std::move(entry));
            }
            crow::json::wvalue result;
            result["entries"] = std::move(arr);
            result["count"]   = records.size();
            return crow::response(result);
        });
}
\end{lstlisting}

\begin{lstlisting}[style=cpp, caption={Server Entry Point: \texttt{src/main.cpp}}]
int main(int argc, char* argv[]) {
    try {
        kvault::EngineConfig config;
        std::string base_dir = "data";

        for (int i = 1; i < argc; ++i) {
            if (std::string(argv[i]) == "--data-dir" && i + 1 < argc)
                base_dir = argv[++i];
        }

        config.wal_directory      = base_dir + "/wal";
        config.sstable_directory  = base_dir + "/sstables";
        config.server_port        = 8080;

        // KVStore constructor: creates directories, initializes subsystems,
        // and replays the WAL to recover from any previous crash.
        auto store = std::make_shared<kvault::KVStore>(config);

        std::cout << "KVault initialized.\n"
                  << "Recovered WAL size:  " << store->wal_size()     << " bytes\n"
                  << "Active SSTables:     " << store->sstable_count() << "\n";

        kvault::ApiServer server(store, config.server_port);
        std::cout << "HTTP API on port " << config.server_port << " -- Ctrl+C to stop.\n";
        server.run(); // Blocks; Crow manages the thread pool internally.

    } catch (const std::exception& e) {
        std::cerr << "Fatal error: " << e.what() << "\n";
        return 1;
    }
    return 0;
}
\end{lstlisting}

\begin{lstlisting}[style=js, caption={React Metrics Hook: \texttt{dashboard/src/hooks/useMetrics.js}}]
import { useState, useEffect, useCallback } from 'react';

const API_BASE = 'http://localhost:8080';

export function useMetrics(pollIntervalMs = 2000) {
    const [metrics, setMetrics] = useState({
        memtable_size_bytes: 0,
        wal_size_bytes:      0,
        sstable_count:       0,
    });
    const [error, setError] = useState(null);

    const fetchMetrics = useCallback(async () => {
        try {
            const res = await fetch(`${API_BASE}/api/metrics`);
            if (!res.ok) throw new Error(`HTTP ${res.status}`);
            const data = await res.json();
            setMetrics(data);
            setError(null);
        } catch (err) {
            setError(err.message);
        }
    }, []);

    useEffect(() => {
        fetchMetrics();
        // Poll /api/metrics every pollIntervalMs to keep the dashboard live.
        const id = setInterval(fetchMetrics, pollIntervalMs);
        return () => clearInterval(id); // Cleanup on unmount.
    }, [fetchMetrics, pollIntervalMs]);

    return { metrics, error };
}
\end{lstlisting}

\section{Deployment: Render Docker Containerization}

\subsection{Conceptual Overview}
To deploy the C++ database engine to cloud infrastructure, containerization is essential to isolate dependencies and ensure a consistent execution environment. Render.com provides a free Docker hosting platform (PaaS) that automatically builds and deploys applications directly from GitHub, operating within an ephemeral storage constraint.

\subsection{The Building Process}
A \texttt{Dockerfile} was created to define a self-contained Ubuntu 22.04 environment. The C++ backend utilizes the standard port \texttt{8080}, which Render automatically detects and routes external HTTP traffic to.

\subsection{Implementation Details}
The \texttt{Dockerfile} installs the required C++20 toolchain (\texttt{g++-13}), \texttt{cmake}, and \texttt{ninja-build}. It copies the source, executes the build pipeline via \texttt{cmake -B build -G Ninja -DCMAKE\_BUILD\_TYPE=Release}, exposes port 8080, and defines the compiled \texttt{kvault\_server} binary as the container entrypoint. 

\begin{lstlisting}[style=cpp, caption={Port configuration update in \texttt{main.cpp}}]
config.wal_directory = base_dir + "/wal";
config.sstable_directory = base_dir + "/sstables";
config.server_port = 7860; // Modified for Hugging Face Spaces compatibility
\end{lstlisting}

\end{document}
